problem:  
Let \( X \) be a smooth \(n\)-dimensional Fano manifold with reductive automorphism group \(\mathrm{Aut}(X)\), and assume that \(X\) admits a unique Kähler–Einstein metric \( \omega_{\mathrm{KE}} \).  

Denote by \( \mathrm{Def}(X) \) the Kuranishi space of complex deformations of \(X\) (i.e., a germ whose tangent space at the origin is \(H^{1}(X, T_X)\) with obstruction potentially in \(H^{2}(X, T_X)\)), and let \( \mathrm{Aut}(X) \) act via its natural holomorphic representation. The GIT quotient \( [\mathrm{Def}(X)//\mathrm{Aut}(X)] \) describes the **algebraic** local moduli of \( \mathbb{Q}\)-Fano varieties near \( X \).  

On the analytic side, let \( \mathcal{M}_{\mathrm{KE}} \) denote the moduli space of Kähler–Einstein Fano manifolds equipped with their metrics, modulo the action of diffeomorphisms isotopic to the identity.  Around \([X,\omega_{\mathrm{KE}}]\), one can consider the local deformation space realized by perturbations of the complex structure and Kähler–Einstein metric solving the coupled Maurer–Cartan and Monge–Ampère equations, modulo automorphisms; denote this germ by \( \mathfrak{M}^{\mathrm{loc}}_{\mathrm{KE}}(X) \).  Topologically, this germ parametrizes nearby complex structures for which a KE metric persists.

**Question:**  
Does the natural mapping  
\[
\Phi: \mathfrak{M}^{\mathrm{loc}}_{\mathrm{KE}}(X) \longrightarrow [\mathrm{Def}(X)//\mathrm{Aut}(X)]
\]
induced by taking the underlying complex structure define a homeomorphism of germs, or even an isomorphism of real-analytic orbifolds?  

Equivalently, is the local Gromov–Hausdorff moduli of Kähler–Einstein Fano manifolds near \(X\) entirely determined (up to real-analytic equivalence) by the complex deformation space modulo \(\mathrm{Aut}(X)\)?  

If not, can one construct an example where the GH–moduli admits additional branches—analytic deformations of the metric that do not originate from any algebraic deformation of the complex structure?

Why is it a "good" problem:  
This problem probes the **local equivalence** between the analytic moduli space of Kähler–Einstein metrics (viewed through deformation theory on the metric side) and the algebraic moduli space governed by \(K\)-stability. Its refined statement replaces informal “metric germ” language with the well-founded Kuranishi and local metric moduli constructions, ensuring mathematical precision.  

Conceptually, it asks: *Are analytic and algebraic moduli locally identical objects?*—a question whose resolution would unify geometric analysis and algebraic geometry at the infinitesimal deformation level.  

If the answer is affirmative, it would yield a local analytic realization of the Yau–Tian–Donaldson correspondence, proving that \(K\)-stability fully controls not only existence but also **deformation** of KE metrics. If false, it would reveal hidden analytic moduli directions invisible to algebraic deformation theory, suggesting new analytic phenomena in the geometry of Fano manifolds.  

The statement is compact and accessible in form yet profoundly deep: it connects complex-analytic deformation theory, geometric analysis of the Monge–Ampère equation, and geometric invariant theory—all within a local framework—thereby bridging major structural themes of modern differential and algebraic geometry.